% !TeX program = lualatex
% ===================================================================
% mwe_book_tagging.tex — teste de viabilidade: book puro + tagging
%
% Objetivo: verificar se \documentclass{book} produz árvore de
% tags correta com latex-lab phase-III, isolando a classe de
% qualquer outro fator (fontes, pacotes editoriais, etc.).
%
% Critério de sucesso: no Acrobat Reader → Tags, a árvore deve
% mostrar <Document> → <Part> → <Sect> → <H1>título de capítulo
% (e não <Part> → <P> como o scrbook está produzindo).
%
% Compilar com: lualatex mwe_book_tagging.tex (duas vezes)
% ===================================================================

\DocumentMetadata{
    lang        = pt-BR,
    pdfstandard = ua-2,
    pdfversion  = 2.0,
    testphase   = {phase-III,table,firstaid},
    tagging     = on,
}

\documentclass[11pt,a5paper,twoside,openright]{book}

\usepackage[brazilian]{babel}
\usepackage{hyperref}
\hypersetup{
    pdftitle           = {Teste book + tagging},
    pdfauthor          = {Teste},
    pdflang            = {pt-BR},
    pdfdisplaydoctitle = true,
}

\title{Teste book + tagging}
\author{Teste}

\begin{document}

\frontmatter

\begin{titlepage}
    \centering
    \vspace*{\fill}
    {\huge\bfseries Teste book + tagging\par}
    \vspace{2\baselineskip}
    {\large Verificação de viabilidade\par}
    \vfill
\end{titlepage}

\tableofcontents

\mainmatter

\part{Primeira parte}

\chapter{Primeiro capítulo}

Texto corrido do primeiro capítulo. Este parágrafo deve virar
uma tag \texttt{<P>} dentro de uma \texttt{<Sect>} que representa
o capítulo, e o título do capítulo acima deve virar \texttt{<H1>}.

\section{Uma seção}

Texto da seção. Deve aparecer como \texttt{<P>} dentro de uma
\texttt{<Sect>} aninhada, com \texttt{<H2>} para o título acima.

\subsection{Uma subseção}

Texto da subseção. Deve produzir \texttt{<H3>} + \texttt{<P>}.

\chapter{Segundo capítulo}

Outro capítulo para verificar que a hierarquia se repete
corretamente. Deve haver uma nova \texttt{<Sect>} irmã da
anterior, não aninhada.

\section{Outra seção}

Mais um teste de \texttt{<H2>}.

\end{document}
